\scanpage{23}\mspage{22}
Il en résulte que $\overline{\{\deleted{\unsure{b}}a\}}=\overline{\{b\}}\Longleftrightarrow a=b$ \unsure{(donc|donc|Alors} l'espace top $\unsure{L|\mathscr{L}}$ satisfait l'axiome $T_0$, il est ``présobre'' i.e. s'injecte \unsure{dans|ds} le présobrifié) -- je ne sais si c'est
le cas pour tout tronçon, \unsure{ni|ni même} si un tronçon, avec sa top. canonique, est sobre (donc, de plus, tte partie fermée \uline{irréductible} aurait un pt générique \dots)
% A step-shaped double vertical stroke joins the underline above to the end of the preceding line.
\margin{\unsure{A|À} voir\unsure{.|,}}
les ``\unsure{segment|segments}'' $[a,b]$ ne sont pas néc. fermés pour la top. canonique\unsure{$^{\dagger}$|$^{+}$|!}\unsure{,|.}
Ce n'est que lorsque $L$ est \uline{tot} ordonné qu'on peut conclure que les pts sont fermés (réciproque ?) et en fait, $L$ est $\unsure{=|\simeq}$ \uline{séparé}, et de plus les $[a,b]$, $L_{\geq a}$, $L_{\leq\deleted{\unsure{b}}a}$ sont fermés.

On peut définir une autre topologie, ayant comme base fondamentale de fermés les $[a,b]$, $L_{\geq a}$, $L_{\leq a}$ (dont on prend les intersections \unsure{quelc\textsuperscript{on}|quelc.} de réunions finies, pour avoir tous les fermés)
